\begin{proof}[Proof of \cref{obj:thm:frontier-answer}]
\begin{enumerate}
\item Fix admissible \(n,d,q\). We first establish the upper bound. Define
\[
p_k(q):=1-(1-q)^k
\quad(k\in\mathbb N,\ q\in[0,1]),
\qquad
p:=p_d(q),
\qquad
c_{\mathrm{exp}}:=1-e^{-1}\in(0,1).
\]
The reveal-probability bounds are
\[
c_{\mathrm{exp}}\min\{1,dq\}\le p\le\min\{1,dq\}.
\]
For the upper bound, \(p\le1\), and induction gives \(p_k(q)\le kq\): the initial value is \(p_0(q)=0\), and
\[
p_{k+1}(q)
=q+(1-q)p_k(q)
\le q+(1-q)kq
\le(k+1)q.
\]
For the lower bound, the exponential inequality \(1-q\le e^{-q}\) yields
\[
(1-q)^d\le e^{-dq}.
\]
When \(dq\le1\), convexity of the exponential gives
\[
e^{-dq}\le(1-dq)+dq\,e^{-1},
\qquad
p\ge c_{\mathrm{exp}}dq.
\]
When \(dq>1\), monotonicity gives
\[
e^{-dq}\le e^{-1},
\qquad
p\ge c_{\mathrm{exp}}.
\]
These prove the claimed reveal-probability bounds; in particular, \(p>0\) when \(q>0\).
% lean: CausalSmith.Experimentation.ThinnedgraphAdditiveRiskFrontier.reveal_probability_ge_chord

For admissible parameters with \(q>0\), define
\[
v_{\mathrm{F}}(n,d,q):=\frac{d^2}{np}.
\]
Since \(p\le1\) and \(p\le dq\),
\[
\frac{d^2}{n}\le v_{\mathrm{F}}(n,d,q),
\qquad
\frac{d}{nq}\le v_{\mathrm{F}}(n,d,q).
\]
Using \(d+1\le2d\) and \(1-q\le1\), the envelope in
\cref{obj:def:upper-envelope} therefore satisfies
\[
u(n,d,q)
=\frac{5(d+1)^2}{n}+\frac{4d(1-q)}{nq}
\le20\frac{d^2}{n}+4\frac{d}{nq}
\le24v_{\mathrm{F}}(n,d,q).
\]
If \(v_{\mathrm{F}}(n,d,q)\le1\), this implies
\[
\min\{1,u(n,d,q)\}\le24v_{\mathrm{F}}(n,d,q)=24F(n,d,q).
\]
If \(v_{\mathrm{F}}(n,d,q)>1\), then
\[
\min\{1,u(n,d,q)\}\le1\le24=24F(n,d,q).
\]
At \(q=0\), the same comparison follows from
\[
u(n,d,0)=+\infty,
\qquad
F(n,d,0)=1.
\]
% lean: CausalSmith.Experimentation.ThinnedgraphAdditiveRiskFrontier.upperEnvelope_le_frontierScale

Use the seeded representative
\[
\widehat T^\star_{n,d,q}(O,\xi):=T^\star_{n,d,q}(O)
\qquad
(O\text{ a complete record},\ \xi\in\mathbb R).
\]
It is measurable by \cref{obj:thm:observable-upper} and
\cref{obj:def:frontier-rule,obj:def:upper-rule}. Those results and the envelope comparison give
\[
R_{n,d}(q)
\le
\sup_{\theta\in\mathcal M_{n,d}}
\mathcal L(\widehat T^\star_{n,d,q};\theta,q)
\le
\min\{1,u(n,d,q)\}
\le24F(n,d,q).
\]
% lean: CausalSmith.Experimentation.ThinnedgraphAdditiveRiskFrontier.frontierEstimator_risk_upper

\item For the lower bound, first suppose that
\[
\frac{d^2}{n}\ge\frac1{16}.
\]
The canonical product design satisfies the assignment, audit, and independence hypotheses of \cref{obj:lem:supplied-block-record-lower}, which supplies
\[
R_{n,d}(q)\ge
\frac1{160000\pi^2}
\min\left\{1,\frac{(d+1)^2}{n}\right\}.
\]
Here
\[
\min\left\{1,\frac{(d+1)^2}{n}\right\}\ge\frac1{16},
\qquad
F(n,d,q)\le1.
\]
Consequently,
\[
R_{n,d}(q)
\ge\frac1{2560000\pi^2}
\ge\frac{F(n,d,q)}{2560000\pi^2}.
\]
% lean: CausalSmith.Experimentation.ThinnedgraphAdditiveRiskFrontier.frontier_minimax_lower

\item Now suppose that
\[
\frac{d^2}{n}<\frac1{16}.
\]
Since \(d\ge1\), we have \(4d\le16d^2<n\). Define the block count and source count by
\[
B:=\left\lfloor\frac{n}{2d}\right\rfloor,
\qquad
m:=Bd.
\]
Floor rounding gives
\[
B\ge1,
\qquad
2m\le n<2m+2d,
\qquad
\frac n4\le m\le\frac n2.
\]
Indeed, \(n>4d\) ensures \(B\ge1\), and \(2d<n/2\) together with
\(n<2m+2d\) gives \(m>n/4\).
% lean: CausalSmith.Experimentation.ThinnedgraphAdditiveRiskFrontier.frontier_block_rounding

Define the testing scale on this parameter range by
\[
s:=s(B,d,q)=
\begin{cases}
\min\{1,d/\sqrt{mp}\},&p>0,\\
1,&p=0.
\end{cases}
\]
Thus \(0<s\le1\). When \(q>0\), we have \(p>0\), and
\[
\left(\frac d{\sqrt{mp}}\right)^2=\frac{d^2}{mp},
\qquad
\frac{d^2}{np}\le\frac{d^2}{mp}.
\]
If \(d/\sqrt{mp}\le1\), these identities imply
\[
F(n,d,q)\le\frac{d^2}{np}\le\frac{d^2}{mp}=s^2.
\]
If \(d/\sqrt{mp}>1\), then \(s=1\) and \(F(n,d,q)\le s^2\).
At \(q=0\), \(p=0\) and \(F(n,d,0)=s^2=1\). Hence, throughout the retention range,
\[
F(n,d,q)\le s^2.
\]
% lean: CausalSmith.Experimentation.ThinnedgraphAdditiveRiskFrontier.frontierScale_le_testingScale_sq

Set
\[
\kappa:=\frac1{100\pi},
\qquad
h:=\kappa s=h_\circ(B,d,q),
\qquad
\Delta:=\frac{mh}{n}.
\]
The positivity of \(m,s,\kappa\), and \(\pi>3\), give
\[
0<h\le\kappa<\frac14,
\qquad
\Delta>0.
\]
Consider the opposite-sign priors of \cref{obj:def:block-prior} at this amplitude. To specify their placement, take the sources to be
\[
\{1,\ldots,m\},
\qquad
C_\ell:=\{m+(\ell-1)d+1,\ldots,m+\ell d\}
\quad(1\le\ell\le B).
\]
The capacity bound \(2m\le n\) places these disjoint recipient blocks within the population. Each recipient has \(d\) incoming arrows, each source has \(d\) outgoing arrows, and the remaining degrees are zero.

The baseline density in \cref{obj:def:block-prior} vanishes outside
\([-1/4,1/4]\), so \(|U_\ell|\le1/4\) almost surely for every block. For either
\(\sigma\in\{-1,1\}\), each recipient \(i\in C_\ell\) consequently satisfies
\[
|a_i|+|t_i|+\sum_{j:(j,i)\in G}|b_{ij}|
=
\left|U_\ell-\frac{\sigma h}{2}\right|
+d\frac hd
\le\frac14+\frac{3h}{2}
\le1.
\]
Source and padding rows have coefficient mass zero. Thus both priors are supported on \(\mathcal M_{n,d}\). Summing the spillover contributions by source gives their constant targets:
\[
\tau(\theta)
=\frac1n\sum_{j=1}^{m}\sum_{i:(j,i)\in G}\frac{\sigma h}{d}
=\frac1n\sum_{j=1}^{m}\sigma h
=\frac{\sigma mh}{n}
=\sigma\Delta.
\]
% lean: CausalSmith.Experimentation.ThinnedgraphAdditiveRiskFrontier.block_support

The record maps are jointly measurable: their graph and assignment coordinates are finite, and their outcome coordinates are affine in the baselines. The canonical product design and \(2Bd\le n\) satisfy the hypotheses of \cref{obj:thm:block-testing-scale}, which supplies
\[
\operatorname{TV}(P_{+1,h},P_{-1,h})\le\frac12.
\]
Applying \cref{obj:lem:full-record-two-prior-risk} to these supported priors yields
\[
R_{n,d}(q)
\ge
\frac{\Delta^2}{2}
\left(1-\operatorname{TV}(P_{+1,h},P_{-1,h})\right)
\ge
\frac{\Delta^2}{4}
=
\frac14\left(\frac{m\kappa s}{n}\right)^2.
\]
This bound applies to the full class of measurable randomized estimators in
\cref{obj:def:risk}.
% lean: CausalSmith.Experimentation.ThinnedgraphAdditiveRiskFrontier.block_testing_minimax_lower

Finally, \(m/n\ge1/4\) gives
\[
\frac{\kappa s}{4}\le\frac{m\kappa s}{n},
\qquad
\frac14\left(\frac{\kappa s}{4}\right)^2
=\frac{s^2}{640000\pi^2}.
\]
Combining the preceding inequalities,
\[
\frac{F(n,d,q)}{2560000\pi^2}
\le
\frac{s^2}{2560000\pi^2}
\le
\frac{s^2}{640000\pi^2}
\le
\frac14\left(\frac{m\kappa s}{n}\right)^2
\le R_{n,d}(q).
\]
Together with the large-degree case and the upper bound, this proves the displayed risk chain.
% lean: CausalSmith.Experimentation.ThinnedgraphAdditiveRiskFrontier.frontier_minimax_lower

\item We next prove necessity in the consistency criterion. Let \(d_n,q_n\) be admissible sequences and suppose that measurable randomized estimators \(T_n\) satisfy
\[
\sup_{\theta\in\mathcal M_{n,d_n}}
\mathcal L(T_n;\theta,q_n)\longrightarrow0.
\]
Define the positive constant
\[
C_{\mathrm{low}}:=2560000\pi^2>0.
\]
For every \(n\ge4\), the lower bound and the definition of minimax risk give
\[
0\le
\frac{F(n,d_n,q_n)}{C_{\mathrm{low}}}
\le R_{n,d_n}(q_n)
\le
\sup_{\theta\in\mathcal M_{n,d_n}}
\mathcal L(T_n;\theta,q_n).
\]
Squeezing in the extended nonnegative reals and passing to real values at zero yields
\[
\frac{F(n,d_n,q_n)}{C_{\mathrm{low}}}\longrightarrow0,
\qquad
F(n,d_n,q_n)\longrightarrow0.
\]

To extract the two ratio limits, we establish the explicit positive-retention comparison. Fix admissible \(n,d\) and \(0<q\le1\). The reveal-probability bounds from the first part imply
\[
c_{\mathrm{exp}}v_{\mathrm{F}}(n,d,q)
\le
\begin{cases}
\dfrac{d}{nq},&dq\le1,\\[4pt]
\dfrac{d^2}{n},&dq>1.
\end{cases}
\]
Indeed, in the first case \(p\ge c_{\mathrm{exp}}dq\), and in the second
\(p\ge c_{\mathrm{exp}}\). Along with the two lower comparisons for
\(v_{\mathrm{F}}\), this gives
\[
\frac{d^2}{n}+\frac{d}{nq}\le2v_{\mathrm{F}}(n,d,q),
\qquad
c_{\mathrm{exp}}v_{\mathrm{F}}(n,d,q)
\le\frac{d^2}{n}+\frac{d}{nq}.
\]
If \(v_{\mathrm{F}}(n,d,q)\ge1\), then
\[
\frac12\min\left\{1,\frac{d^2}{n}+\frac{d}{nq}\right\}
\le1=F(n,d,q).
\]
If \(v_{\mathrm{F}}(n,d,q)<1\), then
\[
\frac12\min\left\{1,\frac{d^2}{n}+\frac{d}{nq}\right\}
\le\frac12\left(\frac{d^2}{n}+\frac{d}{nq}\right)
\le v_{\mathrm{F}}(n,d,q)=F(n,d,q).
\]
For the other direction, \(0<c_{\mathrm{exp}}<1\) and
\(F(n,d,q)=\min\{1,v_{\mathrm{F}}(n,d,q)\}\) imply
\[
c_{\mathrm{exp}}F(n,d,q)\le1,
\qquad
c_{\mathrm{exp}}F(n,d,q)
\le\frac{d^2}{n}+\frac{d}{nq}.
\]
Therefore
\[
\frac12\min\left\{1,\frac{d^2}{n}+\frac{d}{nq}\right\}
\le F(n,d,q)
\le
c_{\mathrm{exp}}^{-1}
\min\left\{1,\frac{d^2}{n}+\frac{d}{nq}\right\}.
\]
% lean: CausalSmith.Experimentation.ThinnedgraphAdditiveRiskFrontier.frontier_elbow

For the sequence argument, define on the tail \(n\ge4\)
\[
\rho_{\mathrm{deg}}(n):=\frac{d_n^2}{n}\in[0,\infty),
\qquad
\rho_{\mathrm{audit}}(n):=
\begin{cases}
+\infty,&q_n=0,\\[2pt]
\dfrac{d_n}{nq_n},&q_n\ne0,
\end{cases}
\quad\in[0,+\infty].
\]
All subsequent sequence bounds concern this tail. Since \(F(n,d_n,q_n)\to0\), eventually
\[
F(n,d_n,q_n)<\frac12.
\]
The identity \(F(n,d_n,0)=1\) forces \(q_n>0\) eventually. The lower comparison then gives
\[
\frac12\min\{1,\rho_{\mathrm{deg}}(n)+\rho_{\mathrm{audit}}(n)\}
\le F(n,d_n,q_n)<\frac12.
\]
Thus the sum is eventually less than one, and
\[
0\le\rho_{\mathrm{deg}}(n)
\le\rho_{\mathrm{deg}}(n)+\rho_{\mathrm{audit}}(n)
\le2F(n,d_n,q_n),
\]
\[
0\le\rho_{\mathrm{audit}}(n)
\le\rho_{\mathrm{deg}}(n)+\rho_{\mathrm{audit}}(n)
\le2F(n,d_n,q_n).
\]
Both ratios converge to zero by squeezing. The audit ratio is eventually finite, so its real convergence also gives convergence in the extended nonnegative reals.
% lean: CausalSmith.Experimentation.ThinnedgraphAdditiveRiskFrontier.precision_frontier_explicit

\item Conversely, suppose that
\[
\rho_{\mathrm{deg}}(n)\longrightarrow0,
\qquad
\rho_{\mathrm{audit}}(n)\longrightarrow0,
\]
with the second limit taken in the extended nonnegative reals. Eventually
\(\rho_{\mathrm{audit}}(n)<1\), which excludes \(q_n=0\); admissibility then gives
\(q_n>0\). Passing to real values at zero yields
\[
\frac{d_n}{nq_n}\longrightarrow0,
\qquad
\rho_{\mathrm{deg}}(n)+\rho_{\mathrm{audit}}(n)\longrightarrow0.
\]
The positive-retention comparison gives, eventually,
\[
0\le F(n,d_n,q_n)
\le c_{\mathrm{exp}}^{-1}
\bigl(\rho_{\mathrm{deg}}(n)+\rho_{\mathrm{audit}}(n)\bigr),
\]
so \(F(n,d_n,q_n)\to0\).
% lean: CausalSmith.Experimentation.ThinnedgraphAdditiveRiskFrontier.frontierScale_tendsto_iff

Choose the measurable seeded attaining estimator for every \(n\):
\[
T_n(O,\xi):=\widehat T^\star_{n,d_n,q_n}(O,\xi).
\]
Its risk satisfies, for every \(n\ge4\),
\[
0\le
\sup_{\theta\in\mathcal M_{n,d_n}}\mathcal L(T_n;\theta,q_n)
\le24F(n,d_n,q_n)\longrightarrow0.
\]
The finitely many indices below four have no effect on the limit. This proves sufficiency and completes the equivalence.
% lean: CausalSmith.Experimentation.ThinnedgraphAdditiveRiskFrontier.precision_frontier_explicit

\item The equality of the total observable rules follows directly from their construction:
\[
T^\star_{n,d,q}(O)=T^{\mathrm{up}}(O)
\qquad
\left(
\begin{array}{l}
V\text{ any finite population},\quad n,d\in\mathbb N,\\
q\in\mathbb R,\quad O\text{ any record on }V
\end{array}
\right).
\]
On the admissible parameter range these are the rules of
\cref{obj:def:frontier-rule,obj:def:upper-rule}, and the measurable seeded representative used above is precisely the representative of the upper rule.
% lean: CausalSmith.Experimentation.ThinnedgraphAdditiveRiskFrontier.frontier_answer

\item Finally, \cref{obj:def:frontier-scale} gives the endpoint identities. At zero retention it gives \(F(n,d,0)=1\); at full retention, \(d\ge1\) gives \(1-(1-1)^d=1\). Hence
\[
F(n,d,0)=1,
\qquad
F(n,d,1)=\min\left\{1,\frac{d^2}{n}\right\}.
\]
The positive-retention comparison established above, with
\(c_{\mathrm{exp}}=1-e^{-1}\), is
\[
\frac12\min\left\{1,\frac{d^2}{n}+\frac{d}{nq}\right\}
\le F(n,d,q)
\le
(1-e^{-1})^{-1}
\min\left\{1,\frac{d^2}{n}+\frac{d}{nq}\right\},
\qquad 0<q\le1.
\]
These are the remaining assertions.
% lean: CausalSmith.Experimentation.ThinnedgraphAdditiveRiskFrontier.frontier_elbow
\end{enumerate}
\end{proof}
